% =============================================================================
% chapters/01_introduction.tex — Introduction
% =============================================================================

\section{Introduction}

This is the first chapter of your thesis. The introduction should provide the necessary background context, state the research problem, outline the scope of the study, and explain the significance of the research topic. It should hook the reader and establish why this study is important.

\subsection{Background}
Describe the broad domain of your research. Provide context for the reader so they understand the setting and the general field of study.

\subsection{Research Problem}
Explain the specific gap in the literature or the practical problem that your research aims to address. Why does this study need to be conducted?

\subsection{Research Objectives and Questions}
State your primary research objectives and the specific research questions or hypotheses that guide this study.

\begin{figure}[htbp]
    \centering
    % Example figure placeholder using LaTeX box
    \fbox{\parbox{0.8\textwidth}{\centering\vspace{3cm}
        \textbf{Figure Placeholder} \\
        \textit{[Insert your project diagram or conceptual framework here]}
        \vspace{3cm}}}
    \caption{Conceptual Framework of the Study}
    \label{fig:conceptual_framework}
\end{figure}

\subsection{Structure of the Thesis}
Briefly outline the contents of the remaining chapters of the thesis: Chapter 2 reviews the relevant literature. Chapter 3 explains the methodology. Chapter 4 presents the results. Chapter 5 discusses the findings, implications, and limitations. Finally, Chapter 6 concludes the study and suggests directions for future research.
